\documentclass[a4paper,11pt,fontset=none]{ctexart}
% ============================================================
%  hjstyle.tex  —  shared preamble for the IMP/HIAF application
%  Use:  \documentclass[a4paper,11pt,fontset=none]{ctexart}
%        % ============================================================
%  hjstyle.tex  —  shared preamble for the IMP/HIAF application
%  Use:  \documentclass[a4paper,11pt,fontset=none]{ctexart}
%        % ============================================================
%  hjstyle.tex  —  shared preamble for the IMP/HIAF application
%  Use:  \documentclass[a4paper,11pt,fontset=none]{ctexart}
%        \input{hjstyle.tex}
%  Compile with: xelatex
% ============================================================

\usepackage{geometry}
\geometry{a4paper, top=2.4cm, bottom=2.3cm, left=2.4cm, right=2.4cm, headsep=10pt}

\usepackage{amsmath,amssymb}
\usepackage{bm}
\usepackage{graphicx}
\usepackage{wrapfig}
\usepackage{float}
\usepackage{array}
\usepackage{booktabs}
\usepackage{tabularx}
\usepackage{enumitem}
\usepackage[dvipsnames,table]{xcolor}
\usepackage{titlesec}
\usepackage{fancyhdr}
\usepackage[normalem]{ulem}
\usepackage{caption}
\usepackage{tikz}
\usetikzlibrary{arrows.meta,positioning,calc,shapes.geometric,backgrounds,fit}
\usepackage{fontspec}

% ---------- Fonts: portable across machines (see README.md) ----------
% Latin = XCharter (bundled with TeX Live). CJK = macOS system fonts when present
% (preserves the original look), else Fandol (bundled with TeX Live) so the document
% builds on ANY TeX Live install (Mac, Linux, Windows) with no extra font setup.
\setmainfont{XCharter}% Latin: XCharter ships with TeX Live (always available)
\IfFontExistsTF{Songti SC}
  {\setCJKmainfont{Songti SC}\newCJKfontfamily\song{Songti SC}}
  {\setCJKmainfont{FandolSong-Regular.otf}[BoldFont=FandolSong-Bold.otf]%
   \newCJKfontfamily\song{FandolSong-Regular.otf}}
\IfFontExistsTF{Heiti SC}
  {\setCJKsansfont{Heiti SC}\newCJKfontfamily\hei{Heiti SC}}% CJK sans for emphasis/headings
  {\setCJKsansfont{FandolHei-Regular.otf}[BoldFont=FandolHei-Bold.otf]%
   \newCJKfontfamily\hei{FandolHei-Regular.otf}}
% route enclosed numbers ①-⑳ and Roman numerals to the CJK font (XCharter lacks them)
\xeCJKDeclareCharClass{CJK}{"2160 -> "2188}
\xeCJKDeclareCharClass{CJK}{"2460 -> "24FF}

% ---------- Colour palette ----------
\definecolor{impblue}{RGB}{31,73,125}     % deep institutional blue
\definecolor{impsky}{RGB}{46,116,181}     % brighter accent blue
\definecolor{impgray}{RGB}{120,120,120}   % header/footer grey
\definecolor{impink}{RGB}{20,40,70}        % near-black ink for key text
\definecolor{impbg}{RGB}{238,243,249}      % pale blue panel background

% ---------- Paragraph style: clean block layout ----------
\setlength{\parindent}{0pt}
\setlength{\parskip}{0.55em}
\linespread{1.06}
\emergencystretch=1.6em   % absorb long unbreakable Latin/math runs, avoid overfull
\setlist{nosep, leftmargin=1.5em, topsep=2pt, itemsep=3pt}

% ---------- Section headings (blue, ruled) ----------
\titleformat{\section}[hang]
  {\Large\bfseries\color{impblue}}{\thesection.}{0.5em}{}
  [{\vspace{1pt}\color{impblue!35}\titlerule[1pt]}]
\titlespacing*{\section}{0pt}{1.1em}{0.5em}

\titleformat{\subsection}[hang]
  {\large\bfseries\color{impsky}}{\thesubsection}{0.5em}{}
\titlespacing*{\subsection}{0pt}{0.8em}{0.25em}

\titleformat{\subsubsection}[runin]
  {\bfseries\color{impink}}{}{0em}{}[\,]
\titlespacing*{\subsubsection}{0pt}{0.4em}{0.5em}

% ---------- Captions ----------
\captionsetup{font=small, labelfont={bf,color=impblue}, labelsep=period, skip=4pt}

% ---------- Running header / footer ----------
\newcommand{\doctitle}{}
\pagestyle{fancy}
\fancyhf{}
\setlength{\headheight}{14pt}
\fancyhead[L]{\small\color{impgray}\,马越\ \textperiodcentered\ Yue Ma}
\fancyhead[R]{\small\color{impgray}\doctitle\,}
\fancyfoot[C]{\small\color{impgray}\thepage}
\renewcommand{\headrulewidth}{0.4pt}
\renewcommand{\footrulewidth}{0pt}
\renewcommand{\headrule}{\hbox to\headwidth{\color{impblue!30}\leaders\hrule height \headrulewidth\hfill}}

% ---------- Physics macros ----------
\newcommand{\hn}[2]{\ensuremath{^{#1}_{\Lambda}\mathrm{#2}}}        % single-Lambda hypernucleus  \hn{3}{H}
\newcommand{\hnn}[2]{\ensuremath{^{#1}_{\Lambda\Lambda}\mathrm{#2}}}% double-Lambda hypernucleus  \hnn{6}{He}
\newcommand{\hsig}[2]{\ensuremath{^{#1}_{\Sigma}\mathrm{#2}}}       % Sigma hypernucleus  \hsig{3}{H}
\newcommand{\dbar}{\ensuremath{\bar{d}}}
\newcommand{\pbar}{\ensuremath{\bar{p}}}
\newcommand{\Kbar}{\ensuremath{\bar{K}}}
\newcommand{\KN}{\ensuremath{\bar{K}N}}
\newcommand{\Kmpp}{\ensuremath{K^{-}pp}}
\newcommand{\KNN}{\ensuremath{\bar{K}NN}}
\newcommand{\Lam}{\ensuremath{\Lambda}}
\newcommand{\Sig}{\ensuremath{\Sigma}}
\newcommand{\Xibar}{\ensuremath{\Xi}}
\newcommand{\Lc}{\ensuremath{\Lambda_c^{+}}}
\newcommand{\gevu}{\,GeV/u}
\newcommand{\mevc}{\,MeV/$c$}

% ---------- Emphasis helpers ----------
\newcommand{\lead}[1]{\textcolor{impblue}{\textbf{#1}}}   % blue bold lead-in
\newcommand{\keypt}[1]{\textbf{#1}}                         % key statement
\newcommand{\keyu}[1]{\uline{\textbf{#1}}}                  % strongest emphasis (use sparingly)
\newcommand{\role}[1]{\textcolor{impsky}{\textbf{[#1]}}}    % authorship/role tag
\newcommand{\tech}[1]{\textmd{#1}}                          % technical term/acronym: keep normal weight even inside bold

% ---------- Blue-bar callout block (for "my contribution" etc.) ----------
\newenvironment{callout}
  {\par\smallskip\noindent\begin{list}{}{%
     \setlength{\leftmargin}{1.1em}\setlength{\rightmargin}{0pt}%
     \setlength{\topsep}{2pt}\setlength{\itemsep}{0pt}}%
     \item[]\color{impink}%
     \def\FrameRule{0pt}}
  {\end{list}\par\smallskip}

% Title banner used at the top of each document
\newcommand{\hjbanner}[3]{% #1 = CN title, #2 = EN subtitle, #3 = meta line
  {\noindent\Huge\bfseries\color{impblue} #1\par}
  \vspace{2pt}
  {\noindent\large\color{impink} #2\par}
  \vspace{4pt}
  {\noindent\small\color{impgray} #3\par}
  \vspace{2pt}
  {\noindent\color{impblue!45}\rule{\linewidth}{1.2pt}\par}
  \vspace{4pt}
}

\usepackage[hidelinks]{hyperref}

%  Compile with: xelatex
% ============================================================

\usepackage{geometry}
\geometry{a4paper, top=2.4cm, bottom=2.3cm, left=2.4cm, right=2.4cm, headsep=10pt}

\usepackage{amsmath,amssymb}
\usepackage{bm}
\usepackage{graphicx}
\usepackage{wrapfig}
\usepackage{float}
\usepackage{array}
\usepackage{booktabs}
\usepackage{tabularx}
\usepackage{enumitem}
\usepackage[dvipsnames,table]{xcolor}
\usepackage{titlesec}
\usepackage{fancyhdr}
\usepackage[normalem]{ulem}
\usepackage{caption}
\usepackage{tikz}
\usetikzlibrary{arrows.meta,positioning,calc,shapes.geometric,backgrounds,fit}
\usepackage{fontspec}

% ---------- Fonts: portable across machines (see README.md) ----------
% Latin = XCharter (bundled with TeX Live). CJK = macOS system fonts when present
% (preserves the original look), else Fandol (bundled with TeX Live) so the document
% builds on ANY TeX Live install (Mac, Linux, Windows) with no extra font setup.
\setmainfont{XCharter}% Latin: XCharter ships with TeX Live (always available)
\IfFontExistsTF{Songti SC}
  {\setCJKmainfont{Songti SC}\newCJKfontfamily\song{Songti SC}}
  {\setCJKmainfont{FandolSong-Regular.otf}[BoldFont=FandolSong-Bold.otf]%
   \newCJKfontfamily\song{FandolSong-Regular.otf}}
\IfFontExistsTF{Heiti SC}
  {\setCJKsansfont{Heiti SC}\newCJKfontfamily\hei{Heiti SC}}% CJK sans for emphasis/headings
  {\setCJKsansfont{FandolHei-Regular.otf}[BoldFont=FandolHei-Bold.otf]%
   \newCJKfontfamily\hei{FandolHei-Regular.otf}}
% route enclosed numbers ①-⑳ and Roman numerals to the CJK font (XCharter lacks them)
\xeCJKDeclareCharClass{CJK}{"2160 -> "2188}
\xeCJKDeclareCharClass{CJK}{"2460 -> "24FF}

% ---------- Colour palette ----------
\definecolor{impblue}{RGB}{31,73,125}     % deep institutional blue
\definecolor{impsky}{RGB}{46,116,181}     % brighter accent blue
\definecolor{impgray}{RGB}{120,120,120}   % header/footer grey
\definecolor{impink}{RGB}{20,40,70}        % near-black ink for key text
\definecolor{impbg}{RGB}{238,243,249}      % pale blue panel background

% ---------- Paragraph style: clean block layout ----------
\setlength{\parindent}{0pt}
\setlength{\parskip}{0.55em}
\linespread{1.06}
\emergencystretch=1.6em   % absorb long unbreakable Latin/math runs, avoid overfull
\setlist{nosep, leftmargin=1.5em, topsep=2pt, itemsep=3pt}

% ---------- Section headings (blue, ruled) ----------
\titleformat{\section}[hang]
  {\Large\bfseries\color{impblue}}{\thesection.}{0.5em}{}
  [{\vspace{1pt}\color{impblue!35}\titlerule[1pt]}]
\titlespacing*{\section}{0pt}{1.1em}{0.5em}

\titleformat{\subsection}[hang]
  {\large\bfseries\color{impsky}}{\thesubsection}{0.5em}{}
\titlespacing*{\subsection}{0pt}{0.8em}{0.25em}

\titleformat{\subsubsection}[runin]
  {\bfseries\color{impink}}{}{0em}{}[\,]
\titlespacing*{\subsubsection}{0pt}{0.4em}{0.5em}

% ---------- Captions ----------
\captionsetup{font=small, labelfont={bf,color=impblue}, labelsep=period, skip=4pt}

% ---------- Running header / footer ----------
\newcommand{\doctitle}{}
\pagestyle{fancy}
\fancyhf{}
\setlength{\headheight}{14pt}
\fancyhead[L]{\small\color{impgray}\,马越\ \textperiodcentered\ Yue Ma}
\fancyhead[R]{\small\color{impgray}\doctitle\,}
\fancyfoot[C]{\small\color{impgray}\thepage}
\renewcommand{\headrulewidth}{0.4pt}
\renewcommand{\footrulewidth}{0pt}
\renewcommand{\headrule}{\hbox to\headwidth{\color{impblue!30}\leaders\hrule height \headrulewidth\hfill}}

% ---------- Physics macros ----------
\newcommand{\hn}[2]{\ensuremath{^{#1}_{\Lambda}\mathrm{#2}}}        % single-Lambda hypernucleus  \hn{3}{H}
\newcommand{\hnn}[2]{\ensuremath{^{#1}_{\Lambda\Lambda}\mathrm{#2}}}% double-Lambda hypernucleus  \hnn{6}{He}
\newcommand{\hsig}[2]{\ensuremath{^{#1}_{\Sigma}\mathrm{#2}}}       % Sigma hypernucleus  \hsig{3}{H}
\newcommand{\dbar}{\ensuremath{\bar{d}}}
\newcommand{\pbar}{\ensuremath{\bar{p}}}
\newcommand{\Kbar}{\ensuremath{\bar{K}}}
\newcommand{\KN}{\ensuremath{\bar{K}N}}
\newcommand{\Kmpp}{\ensuremath{K^{-}pp}}
\newcommand{\KNN}{\ensuremath{\bar{K}NN}}
\newcommand{\Lam}{\ensuremath{\Lambda}}
\newcommand{\Sig}{\ensuremath{\Sigma}}
\newcommand{\Xibar}{\ensuremath{\Xi}}
\newcommand{\Lc}{\ensuremath{\Lambda_c^{+}}}
\newcommand{\gevu}{\,GeV/u}
\newcommand{\mevc}{\,MeV/$c$}

% ---------- Emphasis helpers ----------
\newcommand{\lead}[1]{\textcolor{impblue}{\textbf{#1}}}   % blue bold lead-in
\newcommand{\keypt}[1]{\textbf{#1}}                         % key statement
\newcommand{\keyu}[1]{\uline{\textbf{#1}}}                  % strongest emphasis (use sparingly)
\newcommand{\role}[1]{\textcolor{impsky}{\textbf{[#1]}}}    % authorship/role tag
\newcommand{\tech}[1]{\textmd{#1}}                          % technical term/acronym: keep normal weight even inside bold

% ---------- Blue-bar callout block (for "my contribution" etc.) ----------
\newenvironment{callout}
  {\par\smallskip\noindent\begin{list}{}{%
     \setlength{\leftmargin}{1.1em}\setlength{\rightmargin}{0pt}%
     \setlength{\topsep}{2pt}\setlength{\itemsep}{0pt}}%
     \item[]\color{impink}%
     \def\FrameRule{0pt}}
  {\end{list}\par\smallskip}

% Title banner used at the top of each document
\newcommand{\hjbanner}[3]{% #1 = CN title, #2 = EN subtitle, #3 = meta line
  {\noindent\Huge\bfseries\color{impblue} #1\par}
  \vspace{2pt}
  {\noindent\large\color{impink} #2\par}
  \vspace{4pt}
  {\noindent\small\color{impgray} #3\par}
  \vspace{2pt}
  {\noindent\color{impblue!45}\rule{\linewidth}{1.2pt}\par}
  \vspace{4pt}
}

\usepackage[hidelinks]{hyperref}

%  Compile with: xelatex
% ============================================================

\usepackage{geometry}
\geometry{a4paper, top=2.4cm, bottom=2.3cm, left=2.4cm, right=2.4cm, headsep=10pt}

\usepackage{amsmath,amssymb}
\usepackage{bm}
\usepackage{graphicx}
\usepackage{wrapfig}
\usepackage{float}
\usepackage{array}
\usepackage{booktabs}
\usepackage{tabularx}
\usepackage{enumitem}
\usepackage[dvipsnames,table]{xcolor}
\usepackage{titlesec}
\usepackage{fancyhdr}
\usepackage[normalem]{ulem}
\usepackage{caption}
\usepackage{tikz}
\usetikzlibrary{arrows.meta,positioning,calc,shapes.geometric,backgrounds,fit}
\usepackage{fontspec}

% ---------- Fonts: portable across machines (see README.md) ----------
% Latin = XCharter (bundled with TeX Live). CJK = macOS system fonts when present
% (preserves the original look), else Fandol (bundled with TeX Live) so the document
% builds on ANY TeX Live install (Mac, Linux, Windows) with no extra font setup.
\setmainfont{XCharter}% Latin: XCharter ships with TeX Live (always available)
\IfFontExistsTF{Songti SC}
  {\setCJKmainfont{Songti SC}\newCJKfontfamily\song{Songti SC}}
  {\setCJKmainfont{FandolSong-Regular.otf}[BoldFont=FandolSong-Bold.otf]%
   \newCJKfontfamily\song{FandolSong-Regular.otf}}
\IfFontExistsTF{Heiti SC}
  {\setCJKsansfont{Heiti SC}\newCJKfontfamily\hei{Heiti SC}}% CJK sans for emphasis/headings
  {\setCJKsansfont{FandolHei-Regular.otf}[BoldFont=FandolHei-Bold.otf]%
   \newCJKfontfamily\hei{FandolHei-Regular.otf}}
% route enclosed numbers ①-⑳ and Roman numerals to the CJK font (XCharter lacks them)
\xeCJKDeclareCharClass{CJK}{"2160 -> "2188}
\xeCJKDeclareCharClass{CJK}{"2460 -> "24FF}

% ---------- Colour palette ----------
\definecolor{impblue}{RGB}{31,73,125}     % deep institutional blue
\definecolor{impsky}{RGB}{46,116,181}     % brighter accent blue
\definecolor{impgray}{RGB}{120,120,120}   % header/footer grey
\definecolor{impink}{RGB}{20,40,70}        % near-black ink for key text
\definecolor{impbg}{RGB}{238,243,249}      % pale blue panel background

% ---------- Paragraph style: clean block layout ----------
\setlength{\parindent}{0pt}
\setlength{\parskip}{0.55em}
\linespread{1.06}
\emergencystretch=1.6em   % absorb long unbreakable Latin/math runs, avoid overfull
\setlist{nosep, leftmargin=1.5em, topsep=2pt, itemsep=3pt}

% ---------- Section headings (blue, ruled) ----------
\titleformat{\section}[hang]
  {\Large\bfseries\color{impblue}}{\thesection.}{0.5em}{}
  [{\vspace{1pt}\color{impblue!35}\titlerule[1pt]}]
\titlespacing*{\section}{0pt}{1.1em}{0.5em}

\titleformat{\subsection}[hang]
  {\large\bfseries\color{impsky}}{\thesubsection}{0.5em}{}
\titlespacing*{\subsection}{0pt}{0.8em}{0.25em}

\titleformat{\subsubsection}[runin]
  {\bfseries\color{impink}}{}{0em}{}[\,]
\titlespacing*{\subsubsection}{0pt}{0.4em}{0.5em}

% ---------- Captions ----------
\captionsetup{font=small, labelfont={bf,color=impblue}, labelsep=period, skip=4pt}

% ---------- Running header / footer ----------
\newcommand{\doctitle}{}
\pagestyle{fancy}
\fancyhf{}
\setlength{\headheight}{14pt}
\fancyhead[L]{\small\color{impgray}\,马越\ \textperiodcentered\ Yue Ma}
\fancyhead[R]{\small\color{impgray}\doctitle\,}
\fancyfoot[C]{\small\color{impgray}\thepage}
\renewcommand{\headrulewidth}{0.4pt}
\renewcommand{\footrulewidth}{0pt}
\renewcommand{\headrule}{\hbox to\headwidth{\color{impblue!30}\leaders\hrule height \headrulewidth\hfill}}

% ---------- Physics macros ----------
\newcommand{\hn}[2]{\ensuremath{^{#1}_{\Lambda}\mathrm{#2}}}        % single-Lambda hypernucleus  \hn{3}{H}
\newcommand{\hnn}[2]{\ensuremath{^{#1}_{\Lambda\Lambda}\mathrm{#2}}}% double-Lambda hypernucleus  \hnn{6}{He}
\newcommand{\hsig}[2]{\ensuremath{^{#1}_{\Sigma}\mathrm{#2}}}       % Sigma hypernucleus  \hsig{3}{H}
\newcommand{\dbar}{\ensuremath{\bar{d}}}
\newcommand{\pbar}{\ensuremath{\bar{p}}}
\newcommand{\Kbar}{\ensuremath{\bar{K}}}
\newcommand{\KN}{\ensuremath{\bar{K}N}}
\newcommand{\Kmpp}{\ensuremath{K^{-}pp}}
\newcommand{\KNN}{\ensuremath{\bar{K}NN}}
\newcommand{\Lam}{\ensuremath{\Lambda}}
\newcommand{\Sig}{\ensuremath{\Sigma}}
\newcommand{\Xibar}{\ensuremath{\Xi}}
\newcommand{\Lc}{\ensuremath{\Lambda_c^{+}}}
\newcommand{\gevu}{\,GeV/u}
\newcommand{\mevc}{\,MeV/$c$}

% ---------- Emphasis helpers ----------
\newcommand{\lead}[1]{\textcolor{impblue}{\textbf{#1}}}   % blue bold lead-in
\newcommand{\keypt}[1]{\textbf{#1}}                         % key statement
\newcommand{\keyu}[1]{\uline{\textbf{#1}}}                  % strongest emphasis (use sparingly)
\newcommand{\role}[1]{\textcolor{impsky}{\textbf{[#1]}}}    % authorship/role tag
\newcommand{\tech}[1]{\textmd{#1}}                          % technical term/acronym: keep normal weight even inside bold

% ---------- Blue-bar callout block (for "my contribution" etc.) ----------
\newenvironment{callout}
  {\par\smallskip\noindent\begin{list}{}{%
     \setlength{\leftmargin}{1.1em}\setlength{\rightmargin}{0pt}%
     \setlength{\topsep}{2pt}\setlength{\itemsep}{0pt}}%
     \item[]\color{impink}%
     \def\FrameRule{0pt}}
  {\end{list}\par\smallskip}

% Title banner used at the top of each document
\newcommand{\hjbanner}[3]{% #1 = CN title, #2 = EN subtitle, #3 = meta line
  {\noindent\Huge\bfseries\color{impblue} #1\par}
  \vspace{2pt}
  {\noindent\large\color{impink} #2\par}
  \vspace{4pt}
  {\noindent\small\color{impgray} #3\par}
  \vspace{2pt}
  {\noindent\color{impblue!45}\rule{\linewidth}{1.2pt}\par}
  \vspace{4pt}
}

\usepackage[hidelinks]{hyperref}

\renewcommand{\doctitle}{简历 · Curriculum Vitae}

% date | content row
\newcommand{\cvrow}[2]{%
  \noindent\makebox[2.7cm][l]{\textcolor{impblue}{\small\bfseries #1}}%
  \parbox[t]{\dimexpr\linewidth-2.8cm\relax}{#2}\par\vspace{3pt}}

\begin{document}

\hjbanner{简历 \quad\normalsize Curriculum Vitae}
  {马越 Yue Ma —— 实验核物理学家 · Experimental Nuclear Physicist}
  {日本理化学研究所（RIKEN）少体系统物理实验室 高级研究员（永久职位）\quad\textbar\quad 出生 1979-07-19\quad\textbar\quad \href{mailto:y.ma@riken.jp}{y.ma@riken.jp} \quad\textbar\quad \href{https://ymaphys.github.io}{ymaphys.github.io} \quad\textbar\quad \href{https://orcid.org/0000-0001-8412-8308}{ORCID 0000-0001-8412-8308}}

% ----------------------------------------------------------------------
\section{个人简介}

实验核物理学家，日本理化学研究所（RIKEN）高级研究员（永久职位），在日本 J-PARC、SPring-8、KEK 与德国 GSI、DESY 等国际大型装置拥有 15 年以上一线研究与领导经验，长期从事青年人才培养与国际合作网络建设。\keypt{现担任 \tech{J-PARC E73} 超氚核寿命 实验发言人，同时担任 \tech{SPring-8 LEPS2} 核子短程关联与 \tech{J-PARC} 反氘核两项实验的提案发言人}。

% ----------------------------------------------------------------------
\section{学术领导与学术任职}
\begin{itemize}
  \item \keypt{\tech{J-PARC E73} 超氚核寿命 实验发言人}——实验已完成取束，\hn{4}{H} 寿命与 \hn{3,4}{H} 产生截面已发表，\hn{3}{H} 寿命分析正在推进。
  \item \keypt{\tech{SPring-8 LEPS2} 核子短程关联 提案发言人}——以 $\Lambda$ 超子为探针研究原子核内部核子间短程关联（在研）。
  \item \keypt{\tech{J-PARC} 反氘核实验提案发言人}——反氘核--原子核相互作用与湮灭机制（在研）。
\end{itemize}

\section{标志性研究成果}
\begin{enumerate}
  \item \keypt{\Kmpp（\KNN）反$K$介子--核束缚态的发现}——作为创始成员，研制多击 TDC 并以偏波分析提取信号，给出国际上首个深束缚 \Kmpp 证据（J-PARC E15）。\,{\small Phys.\ Lett.\ B \textbf{789} (2019); Phys.\ Rev.\ C \textbf{102} (2020).}
  \item \keypt{\hn{4}{H} 超核寿命的精密测量}——作为发言人兼通讯作者，提出单 $\gamma$ 标定快 $\pi^0$ 的关键思想并主导测量，$\tau=206\pm8\pm12$\,ps（J-PARC E73）。\,{\small Phys.\ Lett.\ B \textbf{845} (2023); \textbf{873} (2026).}
  \item \keypt{\KN($I{=}1$) 阈附近 $\Lambda\pi$ 结构的首次观测}——作为第一作者主导 Belle 数据分析。\,{\small Phys.\ Rev.\ Lett.\ \textbf{130} (2023).}
  \item \keypt{\tech{TES} 微量能器在强子原子能谱的首次应用}——实现 SDD 数据分析与能量标定，亚电子伏精度（J-PARC E62）。\,{\small Prog.\ Theor.\ Exp.\ Phys.\ \textbf{2016}, 091D01; Phys.\ Rev.\ Lett.\ \textbf{128} (2022).}
\end{enumerate}


% ----------------------------------------------------------------------
\begin{minipage}[t]{0.49\linewidth}
\section{教育背景}
\cvrow{2006--2009}{\textbf{理学博士}，日本东北大学（仙台）}
\cvrow{2004--2006}{理学硕士，日本东北大学（仙台）}
\cvrow{1998--2002}{\textbf{理学学士}（核物理），兰州大学}
\end{minipage}\hfill
\begin{minipage}[t]{0.49\linewidth}
\section{工作经历}
\cvrow{2012--至今}{\textbf{高级研究员（永久职位）}，RIKEN}
\cvrow{2009--2012}{博士后，德国美因茨亥姆霍兹研究所}
\cvrow{2002--2004}{研究实习员，中国科学院近代物理研究所}
\end{minipage}

% ----------------------------------------------------------------------
\begin{minipage}[t]{0.52\linewidth}
\section{科研基金}
\cvrow{2026--2027}{学術変革領域研究（A）公募研究 · 在 J-PARC 探索 Σ 超氚核（\hsig{3}{H}）束缚态（肥山詠美子领域代表「精密数値計算が切り拓く宇宙の量子物質科学」）。\textbf{负责人}}
\cvrow{2024--2028}{基盘研究（A）· 光子束测核内 $\rho$ 介子质量变化。共同研究者}
\cvrow{2023--2025}{RIKEN 奖励课题 · 探索物质--反物质相互作用。\textbf{负责人}}
\cvrow{2017--2020}{JSPS 若手研究（A）17H04842 · 解决超氚寿命之谜。\textbf{负责人}}
\end{minipage}\hfill
\begin{minipage}[t]{0.44\linewidth}
\section{学生指导}
\cvrow{2017--2020}{C.\ Han（博士），中国科学院大学}
\cvrow{2012--2014}{Q.\ Zhang（博士），兰州大学}
\cvrow{2010--2011}{D.\ Lin（博士）；Christina、Pascal（硕士），美因茨大学}
\end{minipage}

% ----------------------------------------------------------------------
\begin{minipage}[t]{0.52\linewidth}
\section{组织与管理 · 特邀报告}
\cvrow{2019--2021}{RIKEN 科研人员大会指导委员会委员；激励研究基金物理学部协调人 / 总协调人}
\cvrow{2022}{特邀报告：第 14 届国际超核与奇异粒子物理大会（HYP 2022，布拉格）}
\cvrow{2025}{特邀报告：第 15 届国际超核与奇异粒子物理大会（HYP 2025，东京）}
\end{minipage}\hfill
\begin{minipage}[t]{0.44\linewidth}
\section{教学 · 语言}
\cvrow{2018/23/24}{强子物理特邀讲师，中国科学院大学（暑期课程）}
\cvrow{语言}{中文（母语）· 英语（流利）· 日语（专业工作水平）}
\end{minipage}

% ----------------------------------------------------------------------
\section{代表性论文 \normalfont\small（Selected Publications，按主导作用与影响力排序；标注作者身份）}

{\small 发表论文 \keypt{100 余篇}（含 Belle / Belle\,II 合作组署名论文，见文末「合作组文章」一节）；下列为第一/通讯作者或起主导作用的代表性成果。}

{\small\color{impgray} 注：日本强子物理论文多按姓氏字母排序，下列标注\textbf{第一作者}/\textbf{通讯作者}/\textbf{发言人}者，为实际主导作用之明确标识。}

{\small
\begin{enumerate}\setlength{\itemsep}{2pt}
  \item \textbf{Y.~Ma} \textit{et al.}, \textit{First observation of $\Lambda\pi^+$ and $\Lambda\pi^-$ signals near the $\bar{K}N(I{=}1)$ mass threshold in $\Lambda_c^+\!\to\Lambda\pi^+\pi^+\pi^-$ decay}, Phys.\ Rev.\ Lett.\ \textbf{130}, 151903 (2023) [Belle]. \role{第一作者}
  \item T.~Akaishi \textit{et al.}, \textit{Precise lifetime measurement of $^{4}_{\Lambda}$H hypernucleus using a novel production method}, Phys.\ Lett.\ B \textbf{845}, 138128 (2023). \role{发言人 \& 通讯作者}
  \item T.~Akaishi \textit{et al.}, \textit{Measurement of $^{3,4}$He$(K^-,\pi^0)$ $^{3,4}_{\Lambda}$H reaction cross section and evaluation of hypertriton $\Lambda$ binding energy}, Phys.\ Lett.\ B \textbf{873}, 140163 (2026). \role{发言人}
  \item \textbf{Y.~Ma} \textit{et al.}, \textit{$\gamma$-ray spectroscopy study of $^{11}_{\Lambda}$B and $^{12}_{\Lambda}$C}, Eur.\ Phys.\ J.\ A \textbf{33}, 243 (2007). \role{第一作者}
  \item S.~Ajimura \textit{et al.}, \textit{$K^-pp$, a $\bar{K}$-meson nuclear bound state, observed in $^3$He$(K^-,\Lambda p)n$ reactions}, Phys.\ Lett.\ B \textbf{789}, 620 (2019). \role{主要作者之一}
  \item T.~Yamaga \textit{et al.}, \textit{Observation of a $\bar{K}NN$ bound state in the $^3$He$(K^-,\Lambda p)n$ reaction}, Phys.\ Rev.\ C \textbf{102}, 044002 (2020). \role{主要贡献者之一}
  \item T.~Hashimoto \textit{et al.}, \textit{Strong-interaction effects in kaonic-helium isotopes at sub-eV precision with X-ray microcalorimeters}, Phys.\ Rev.\ Lett.\ \textbf{128}, 112503 (2022). \role{主要贡献者之一}
  \item B.~S.~Henderson \textit{et al.}, \textit{Hard two-photon contribution to elastic lepton-proton scattering (OLYMPUS)}, Phys.\ Rev.\ Lett.\ \textbf{118}, 092501 (2017). \role{主要贡献者之一；主导亮度监测器（PbF$_2$）研制}  \item K.~Hosomi \textit{et al.}, \textit{Precise determination of $^{12}_{\Lambda}$C level structure by $\gamma$-ray spectroscopy}, Prog.\ Theor.\ Exp.\ Phys.\ \textbf{2015}, 081D01 (2015). \role{第二作者}
\end{enumerate}
}

% ----------------------------------------------------------------------
\section{合作组文章 \normalfont\small（Belle / Belle\,II Collaboration Papers）}

{\small 作为 Belle / Belle\,II 合作组正式成员，自 \keypt{2023 年}起累计署名合作组论文 \keypt{70 余篇}；按研究主题归类，各列代表性论文如下。}

{\small\setlength{\emergencystretch}{3em}

\smallskip
\noindent\lead{一、$B$ 介子 CP 破坏与 CKM 矩阵元精测}\quad{\color{impgray}时间依赖 CP 破坏（$\sin 2\phi_1$ 等）、$|V_{cb}|$/$|V_{ub}|$ 与 CKM 角 $\phi_3$ 的精确测定。}
\begin{enumerate}\setlength{\itemsep}{1.5pt}
  \item \textit{Measurement of $\sin 2\phi_{1}$ in $B^{0}\!\to J/\psi K_{S}^{0}$ decays with a graph-neural-network flavor tagger}, Phys.\ Rev.\ D \textbf{110}, 012001 (2024)\;[Belle\,II].
  \item \textit{Measurement of CP violation in $B^{0}\!\to K_{S}^{0}\pi^{0}$ decays}, Phys.\ Rev.\ Lett.\ \textbf{131}, 111803 (2023)\;[Belle\,II].
  \item \textit{Determination of $|V_{cb}|$ using $\bar{B}^{0}\!\to D^{*+}\ell^{-}\bar\nu_{\ell}$ decays}, Phys.\ Rev.\ D \textbf{108}, 092013 (2023)\;[Belle\,II].
  \item \textit{Measurement of inclusive $B\!\to X_{u}\ell\nu$ partial branching fractions and $|V_{ub}|$}, Phys.\ Rev.\ D \textbf{113}, 032004 (2026)\;[Belle\,II].
  \item \textit{Determination of the CKM angle $\phi_{3}$ from a combination of Belle and Belle\,II results}, J.\ High Energy Phys.\ \textbf{10} (2024) 143 [Belle/Belle\,II].
\end{enumerate}

\smallskip
\noindent\lead{二、稀有衰变与轻子普适性（新物理探针）}\quad{\color{impgray}$B^+\!\to K^+\nu\bar\nu$ 的\keypt{首个实验证据}、$R(D^{*})$ 等轻子普适性检验与 $\tau$ 轻子味破坏（LFV）搜寻。}
\begin{enumerate}\setlength{\itemsep}{1.5pt}
  \item \textit{Evidence for $B^{+}\!\to K^{+}\nu\bar\nu$ decays}, Phys.\ Rev.\ D \textbf{109}, 112006 (2024)\;[Belle\,II].
  \item \textit{Test of light-lepton universality with a measurement of $R(D^{*})$ using hadronic $B$ tagging}, Phys.\ Rev.\ D \textbf{110}, 072020 (2024)\;[Belle\,II].
  \item \textit{First measurement of $R(X_{\tau/\ell})$ as an inclusive test of the $b\!\to c\tau\nu$ anomaly}, Phys.\ Rev.\ Lett.\ \textbf{132}, 211804 (2024)\;[Belle\,II].
  \item \textit{Tests of light-lepton universality in angular asymmetries of $B^{0}\!\to D^{*-}\ell\nu$ decays}, Phys.\ Rev.\ Lett.\ \textbf{131}, 181801 (2023)\;[Belle\,II].
  \item \textit{Measurement of the $B^{+}\!\to\tau^{+}\nu_{\tau}$ branching fraction with a hadronic tagging method}, Phys.\ Rev.\ D \textbf{112}, 072002 (2025)\;[Belle\,II].
  \item \textit{First search for $B\!\to X_{s}\nu\bar\nu$ decays}, Phys.\ Rev.\ Lett.\ \textbf{136}, 231801 (2026)\;[Belle\,II].
  \item \textit{Search for lepton-flavor-violating $\tau$ decays to $\ell\alpha$}, J.\ High Energy Phys.\ \textbf{08} (2025) 155 [Belle].
  \item \textit{Search for a $\tau^{+}\tau^{-}$ resonance in $e^{+}e^{-}\!\to\mu^{+}\mu^{-}\tau^{+}\tau^{-}$ events}, Phys.\ Rev.\ Lett.\ \textbf{131}, 121802 (2023)\;[Belle\,II].
\end{enumerate}

\smallskip
\noindent\lead{三、粲/奇异重子谱学与奇异强子}{\color{impgray}（与本人研究方向高度契合）}\quad{\color{impgray}$\Xi_c$ 衰变与 CP、$D_{s0}^{*}(2317)$ 辐射衰变\keypt{首次观测}、$P_{cs}$ 五夸克态与 \keypt{$\Xi^0 p$/$\Omega^- p$/$\Omega^- n$ 双重子}搜寻。}
\begin{enumerate}\setlength{\itemsep}{1.5pt}
  \item \textit{Observation of the radiative decay $D_{s0}^{*}(2317)^{+}\!\to D_{s}^{*+}\gamma$}, Phys.\ Rev.\ Lett.\ \textbf{136}, 241901 (2026) [Belle/Belle\,II].
  \item \textit{Search for $P_{cs}(4459)$ and $P_{cs}(4338)$ in $\Upsilon(1S,2S)$ inclusive decays}, Phys.\ Rev.\ Lett.\ \textbf{135}, 041901 (2025) [Belle/Belle\,II].
  \item \textit{Search for CP violation in $\Xi_{c}^{+}\!\to\Sigma^{+}h^{+}h^{-}$ and $\Lambda_{c}^{+}\!\to p\,h^{+}h^{-}$ decays}, Phys.\ Rev.\ D \textbf{113}, 032017 (2026)\;[Belle\,II].
  \item \textit{Observation of the singly Cabibbo-suppressed decays $\Xi_{c}^{+}\!\to pK_{S}^{0}$, $\Lambda\pi^{+}$, $\Sigma^{0}\pi^{+}$}, J.\ High Energy Phys.\ \textbf{03} (2025) 061 [Belle/Belle\,II].
  \item \textit{Measurement of branching fractions of $\Xi_{c}^{0}\!\to\Xi^{0}\pi^{0},\,\Xi^{0}\eta,\,\Xi^{0}\eta'$}, J.\ High Energy Phys.\ \textbf{10} (2024) 045 [Belle/Belle\,II].
  \item \textit{Observation of $B^{+}\!\to\Sigma_{c}(2455)^{++}\bar\Xi_{c}^{-}$ and $B^{0}\!\to\Sigma_{c}(2455)^{0}\bar\Xi_{c}^{0}$}, Phys.\ Rev.\ D \textbf{112}, L051101 (2025)\;[Belle/Belle\,II].
\end{enumerate}
}
\end{document}
